%% glossary.tex — Alphabetical technical glossary
%% JA4proxy Technical Reference Manual
%% Note: \chapter{Glossary} is in the main file. This file provides only the entries.

\index{glossary}

\newthought{Terms are listed alphabetically.} Cross-references point to the chapter or
section where each concept is covered in detail.

\begin{description}[leftmargin=0pt, style=nextline, itemsep=8pt]

  \item[Action decider\index{action decider}]
  The final pipeline layer (Layer 14) that maps a computed risk score and the current
  dial setting to one of the six actions: allow, flag, rate\_limit, tarpit, block, or
  ban. The action decider applies the operator's dial policy to the objective score
  produced by the risk scorer.
  See Chapter~\ref{ch:pipeline}, Section on the Action Decider.

  \item[ALPN (Application-Layer Protocol Negotiation)\index{ALPN}]
  A TLS extension (extension type 0x0010) in the ClientHello that lists the
  application-layer protocols the client supports, in preference order. The most common
  values for HTTPS are \code{h2} (HTTP/2) and \code{h1} (HTTP/1.1). \ja\ uses ALPN
  to identify browser traffic — real browsers always negotiate h2 or h1.
  See Chapter~\ref{ch:pipeline}, Layer 1.

  \item[ASN (Autonomous System Number)\index{ASN}]
  A globally unique number assigned to a network operator (ISP, cloud provider,
  university, etc.) by a Regional Internet Registry (RIR). IP prefixes are announced
  under an ASN in the global BGP routing table. \ja\ uses the MaxMind GeoLite2 ASN
  database to map client IPs to their originating ASN and organisation name.
  See Chapter~\ref{ch:signals}, Signal 4.

  \item[AbuseIPDB\index{AbuseIPDB}]
  A community-maintained reputation database of IP addresses reported for malicious
  activity (scanning, brute-force, spam, C2 traffic). Each IP receives a confidence
  score 0–100. \ja\ caches AbuseIPDB responses in Redis and uses the score as a
  weighted signal. Scores below 50 are suppressed to avoid false positives from
  shared infrastructure.
  See Chapter~\ref{ch:signals}, Signal 7.

  \item[Attacker attribution\index{attacker attribution}]
  Pipeline Layer 11. The process of correlating JA4, JA4X, and JA4T fingerprints
  across multiple source IPs to identify coordinated campaigns, botnet activations,
  and IP-rotation evasion attempts. Attribution findings are written by the analytics
  node and read back by the pipeline.
  See Chapter~\ref{ch:pipeline}, Section on Layer 11.

  \item[Asymmetry principle\index{asymmetry principle}]
  The core design principle of \ja: the cost of a false positive (real user blocked)
  is higher than the cost of a false negative (bad bot allowed). This principle
  governs all threshold defaults, caching TTLs, and fallback behaviours. When in
  doubt, fail open.
  See Chapter~\ref{ch:architecture}, Section on Design Principles.

  \item[Ban TTL\index{ban TTL}]
  The time-to-live of an IP ban entry in Redis (\code{ban:\{ip\}}). Initial TTL is
  configurable (default 300 seconds); it escalates by doubling on each re-ban up
  to a maximum of 604800 seconds (7 days). After the TTL expires, the ban entry is
  automatically removed by Redis and the IP is re-evaluated on next connection.
  See Chapter~\ref{ch:pipeline}, Section on Blocking Escalation.

  \item[Beaconing\index{beaconing}]
  A network traffic pattern characterised by regular, periodic connections with low
  variance in inter-arrival times (IAT). Beaconing is characteristic of C2 (command
  and control) implants and automated scheduled tasks, as opposed to human browsing
  which has highly variable timing. Detected using the coefficient of variation of
  IAT.
  See Chapter~\ref{ch:signals}, Signal 6.

  \item[Bloom filter\index{Bloom filter}]
  A probabilistic data structure that supports membership queries with a controlled
  false positive rate and no false negatives. \ja\ uses Redis Bloom filters (via
  RedisBloom module) for enrichment deduplication — to avoid re-querying AbuseIPDB
  or RDAP for IPs that have already been processed. If RedisBloom is unavailable,
  a Redis SET with a 24h TTL is used as fallback.
  See Chapter~\ref{ch:redis}.

  \item[Bypass condition\index{bypass condition}]
  A fast-path rule that short-circuits the pipeline immediately when a specific
  condition is met, without running the signal collection or scoring layers.
  \ja\ has eight bypass conditions, each independently configurable.
  See Chapter~\ref{ch:pipeline}, Section on the Fast-Path Bypass System.

  \item[CIDR (Classless Inter-Domain Routing)\index{CIDR}]
  IP prefix notation combining a base address with a prefix length (e.g.\
  \code{185.220.96.0/21}). Used in \ja\ for subnet blocking, trusted upstream CIDR
  verification, and block expansion. CIDR matching is always performed in-process
  using a trie data structure — never via Redis queries.
  See Chapter~\ref{ch:redis}.

  \item[ClientHello\index{ClientHello}]
  The first message sent by a TLS client in a TLS handshake. It is transmitted in
  plaintext (TLS record type 0x16, content type 0x01) before any encryption is
  established. It contains the TLS version, cipher suite list, extensions (including
  SNI, ALPN, supported groups), and other parameters. \ja\ reads and parses the
  ClientHello without participating in the TLS handshake.
  See Chapter~\ref{ch:architecture}, Section on Data Flow.

  \item[Cold cache\index{cold cache}]
  The state of the local in-process LRU cache immediately after startup, or for a
  newly seen IP address. Cold cache state means there is no locally cached decision;
  the pipeline must execute in full (including Redis lookups and signal collection)
  to produce a decision for this IP.
  Contrast with \textit{warm cache}.

  \item[Coefficient of variation (CV)\index{coefficient of variation}]
  A statistical measure of dispersion: CV = standard deviation / mean. Used in the
  beaconing detector to quantify the regularity of inter-arrival times. Low CV
  ($<$0.4) indicates regular, beaconing-like timing. High CV ($>$0.8) indicates
  variable, human-like timing.
  See Chapter~\ref{ch:signals}, Signal 6.

  \item[Dial\index{dial}]
  A scalar integer setting (0–100) that controls how aggressively risk scores are
  translated into blocking actions. At dial=0, the proxy is in monitor-only mode —
  all connections are allowed regardless of score. At dial=100, configured thresholds
  apply fully. The dial should be raised incrementally after validating the score
  distribution.
  See Chapter~\ref{ch:configuration}, Section on dial.

  \item[DGA (Domain Generation Algorithm)\index{DGA}]
  An algorithm used by malware to generate pseudo-random domain names for C2
  communication, making blacklisting difficult. DGA-generated domains typically
  have high Shannon entropy and use unusual TLD patterns. \ja\ detects DGA-like
  SNI values as a scored signal.
  See Chapter~\ref{ch:signals}, Signal 2.

  \item[FCrDNS (Forward-confirmed reverse DNS)\index{FCrDNS}]
  A DNS validation technique: the PTR record for an IP resolves to a hostname, and
  that hostname's A/AAAA record resolves back to the original IP. FCrDNS is a
  lightweight indicator of network infrastructure legitimacy — residential ISPs,
  major cloud providers, and CDNs maintain consistent FCrDNS.
  See Chapter~\ref{ch:signals}, Signal 5.

  \item[Fail open\index{fail open}]
  The principle that in all uncertainty or failure conditions, connections are allowed
  rather than blocked. An unavailable external service, a Redis error, or an
  unrecognised fingerprint all result in a zero signal contribution — the score can
  only decrease when services fail, not increase. The fail-open guarantee is fundamental
  to \ja's false positive minimisation.
  See Chapter~\ref{ch:security}, Section on the Fail-Open Guarantee.

  \item[Fast path\index{fast path}]
  The set of pipeline layers (0–2) that short-circuit immediately without running
  the signal collection or scoring layers. Fast-path decisions (from local cache,
  bypass conditions) complete in under 100 microseconds. Connections that bypass
  the fast path proceed to the signal collection zone.
  See Chapter~\ref{ch:pipeline}.

  \item[HyperLogLog\index{HyperLogLog}]
  A probabilistic data structure for counting distinct elements with fixed memory
  cost (12KB per key in Redis). Used in \ja\ to count unique IPs per /24 (IPv4)
  or /48 (IPv6) CIDR prefix. Standard error is 0.81\%. Individual elements cannot
  be removed from a HyperLogLog without clearing the entire structure.
  See Chapter~\ref{ch:redis}.

  \item[IAT (Inter-arrival time)\index{IAT}]
  The time elapsed between consecutive connections from the same source (IP + JA4
  fingerprint pair). IAT analysis is the core of the beaconing detection algorithm.
  IAT values are stored as timestamps in a Redis Sorted Set with a 30-minute
  sliding window.
  See Chapter~\ref{ch:signals}, Signal 6.

  \item[JA4\index{JA4}]
  A TLS client fingerprinting algorithm that produces a deterministic hash from the
  TLS ClientHello: cipher suites (sorted), TLS extensions (sorted), ALPN values, and
  signature algorithms. Produces a string of the form
  \fingerprint{t13d1516h2\_8daaf6152771\_02713d6af862}. The algorithm was developed
  by FoxIO and is defined at \url{https://github.com/FoxIO-LLC/ja4}.
  See Chapter~\ref{ch:signals}.

  \item[JA4X\index{JA4X}]
  The extended JA4 fingerprint that includes the TLS extension type and ordering
  information — not just which extensions are present, but in what order the client
  sent them. JA4X can distinguish clients that share the same JA4 value but have
  different extension ordering (e.g., a real Chrome vs a JA4-copying tool).
  See Chapter~\ref{ch:signals}.

  \item[JA4T\index{JA4T}]
  The TCP transport fingerprint component of the JA4 family. Encodes the TCP window
  size, TCP options (MSS, window scale, SACK, timestamps), and their ordering.
  JA4T fingerprints the OS-level TCP/IP stack, making it harder to spoof than JA4
  (which fingerprints the TLS library layer).
  See Chapter~\ref{ch:signals}, Signal 3.

  \item[Monitor mode\index{monitor mode}]
  The operating mode when \code{dial: 0}. In monitor mode, the full pipeline runs
  (all signals collected, scores computed, actions decided), but every connection is
  allowed regardless of score. Actions are logged as \badgemonitor. Monitor mode
  is the default and recommended starting point for all deployments.
  See Chapter~\ref{ch:pipeline}, Section on the Action Decider.

  \item[mTLS (Mutual TLS)\index{mTLS}]
  TLS authentication where both the server and the client present certificates.
  A valid mTLS client certificate is cryptographic proof of identity — the client
  holds the private key corresponding to a certificate signed by the configured
  trusted CA (\code{config/trusted\_cas.pem}). By default, mTLS clients bypass
  the scoring pipeline entirely.
  See Chapter~\ref{ch:pipeline}, Layer 1c.

  \item[PROXY protocol\index{PROXY protocol}]
  A binary or text framing format prepended to a TCP stream by a load balancer to
  convey the original client IP address and port. \ja\ requires HAProxy to use
  PROXY protocol v2 (binary format) so that the real client IP is available after
  the TCP connection arrives from HAProxy.
  See Chapter~\ref{ch:architecture}, Section on Stage 1.

  \item[RDAP (Registration Data Access Protocol)\index{RDAP}]
  The modern replacement for WHOIS. Provides structured JSON responses about IP
  address registrations from Regional Internet Registries (RIRs: ARIN, RIPE, APNIC,
  LACNIC, AFRINIC). \ja\ uses RDAP to identify the organisation owning an IP block
  and cross-reference it against the known-bad organisations list.
  See Chapter~\ref{ch:signals}, Signal 8.

  \item[Risk score\index{risk score}]
  A scalar integer in the range 0–100 produced by the risk scorer (pipeline Layer 13)
  by confidence-weighted aggregation of all signal contributions. A score of 0 means
  clean; 100 means maximum risk. The dial setting controls how the score is
  translated into a blocking action.
  See Chapter~\ref{ch:pipeline}, Section on the Risk Scorer.

  \item[Signal\index{signal}]
  A scored observation produced by one of the eight signal checks (pipeline Layers
  3–10). Each signal returns a \code{RiskSignal(score, weight, reason)} object.
  Signals run concurrently via \code{asyncio.gather()}. A signal that cannot produce
  a result (service unavailable) returns score=0 (fail open).
  See Chapter~\ref{ch:signals}.

  \item[SNI (Server Name Indication)\index{SNI}]
  A TLS extension (type 0x0000) in the ClientHello that specifies the hostname the
  client intends to connect to. SNI allows a single server to host multiple TLS
  virtual hosts on the same IP. \ja\ uses SNI to detect missing SNI, IP-literal SNI,
  and DGA-like domain names.
  See Chapter~\ref{ch:signals}, Signal 2.

  \item[Spamhaus DROP/EDROP\index{Spamhaus DROP}\index{Spamhaus EDROP}]
  The Spamhaus Don't Route Or Peer (DROP) and Extended DROP (EDROP) blocklists
  contain IP ranges that have been hijacked or directly allocated to cybercriminals.
  False positive rate is near zero. \ja\ downloads these lists at startup and on a
  refresh interval, storing them in an in-process trie for O(log n) lookup.
  See Chapter~\ref{ch:pipeline}, Layer 0e.

  \item[Tarpit\index{tarpit}]
  An action that accepts a connection but delivers data to the backend at an
  artificially throttled rate, consuming the attacker's connection slots and CPU
  while expending minimal resources on the defender's side. Effective against
  connection-limited scanners and brute-force tools. \ja\ implements tarpit as a
  rate-controlled TCP relay.
  See Chapter~\ref{ch:pipeline}, Section on the Six Actions.

  \item[TLS fingerprinting\index{TLS fingerprinting}]
  The technique of computing a deterministic identifier from the parameters in a
  TLS ClientHello (cipher suites, extensions, ALPN, etc.) to identify the TLS
  client implementation. Different applications produce different fingerprints even
  when connecting to the same server. The JA4 algorithm is the fingerprinting
  standard used by \ja.
  See Chapter~\ref{ch:signals}.

  \item[Warm cache\index{warm cache}]
  The state of the local in-process LRU cache when it contains a recent decision
  for a given IP address. A warm cache hit means the pipeline can return a decision
  in under 100 microseconds without any Redis or signal collection overhead. Cache
  entries have asymmetric TTLs: ALLOW entries are cached for 30 minutes; BLOCK
  entries for 30 seconds.
  Contrast with \textit{cold cache}.

\end{description}
